\subsection{2D Isentropic Vortex}
\label{SEC:Testcases:IsentropicVortex}

\begin{figure}[h!tb]
  \begin{center}
    \insertimage[width=3.0in]{test_cases/euler_vortex}
    \caption{Isentropic Euler vortex test case.}
    \label{FIG:IsentropicVortex}
  \end{center}
\end{figure}

This test case simulates the long term transport of a vortex by a
uniform base flow over a 2D doubly periodic domain,
as shown in \fref{FIG:IsentropicVortex}.
%
It is aimed at testing a high-order method’s capability to preserve
vorticity in an unsteady inviscid flow.

The flow solver will internally generate a uniform grid for this case
if \inp[nmlGrid::GridFormat]{Grid\_Format}=-1.
%
Other parameters for this case are specified through the
\inp[nmlTestcaseVortex]{nml\_testcase\_vortex} namelist block.
%
Grid resolution can be adjusted by changing the
\inp[nmlGrid::InternalGridDoF]{Internal\_Grid\_DoF}
input parameter which specifies the degree-of-freedom (DoF)
to be applied to each direction.
%
The number of grid cells in each direction is determined as:
%
\begin{equation*}
    \left[ \text{\# of cells} \right]
  = \mathrm{nint}
    ( \text{\inp[nmlGrid::InternalGridDoF]{Internal\_Grid\_DoF}}
    / ( \text{\inp[nmlNumerics::SolutionOrder]{Solution\_Order}} + 1 )
    ).
\end{equation*}
%
For example:\\
\begin{tabular}{@{\qquad}l}
  \inptext{Internal\_Grid\_DoF} = 200 \\
  \inptext{Solution\_Order} = 5 \\
  Grid = 33x33 Cells
\end{tabular}

This case reaches its final simulation time in
about 30 seconds using 52 MPI processes.
%
A detailed validation of the GFR flow solver for the isentropic vortex
test case can be found in \linecite{Spiegel_2015}.

\SeeAlso{
\inp[nmlInitialization::Initialization]{Initialization}=5,
\inp[nmlGrid::InternalGridDoF]{Internal\_Grid\_DoF},
\inp[nmlTestcaseVortex]{nml\_testcase\_vortex}
}

